\chapter{Somme des angles d'un triangle}

\begin{itemize}
\it Soient deux triangles $ABC$ et $A'B'C'$ ayant leurs côtés respectivement parallèles. $(B'C')$ coupe les droites $(AB)$ et $(AC)$ en $D$ et $E$. 
\begin{enumerate}
\item Comparer les angles des triangles $ABC$ et $ADE$, puis ceux des triangles $ADE$ et $A'B'C'$. 
\item Énoncer la propriété qui en résulte pour deux triangles qui ont leurs côtés parallèles.
\end{enumerate}
\it Soient deux demi-droites $[Ax)$ et $[By)$ situées d'un même côté de la droite $(AB)$ et $[Ax')$ et $[By')$ les demi-droites opposées. On suppose $\widehat{BAx}+\widehat{ABy}< 180^o$.
\begin{enumerate}
\item Montrer que les droites $(xx')$ et $(yy')$ sont sécantes et que les demi-droites $[Ax')$ et $[By')$ n'ont pas de point commun.
\item En déduire que $[Ax)$ et $[By)$ se coupent en un point $C$ et énoncer la condition pour que deux demi-droites aient un point commun. 
\end{enumerate}
\it Soit un triangle $ABC$, dont on note les mesures des angles $\widehat{A}$, $\widehat{B}$ et $\widehat{C}$.\begin{enumerate}
\item Évaluer en fonction de $\widehat{B}$ et $\widehat{C}$ l'angle de la hauteur $(AA')$ issue de $A$ et de la  bissectrice intérieure $(AD)$ de l'angle au sommet $A$.
\item Évaluer en fonction de $A$, l'angle $\widehat{BIC}$ des bissectrices intérieures aux sommets $B$ et $C$. 
\item Évaluer en fonction de $\widehat{A}$, l'angle $\widehat{BHC}$ des hauteurs issues de $B$ et $C$. \\
Application numérique : $\widehat{B} = 57^o$, $\widehat{C} = 75^o$.
\end{enumerate}
\it Soit un triangle $ABC$ dans lequel $AC>AB$. On mène la bissectrice de l'angle en $A$ qui coupe $[BC]$ en $D$, puis la perpendiculaire $(BE)$ à $(AD)$. Évaluer en fonction des angles en $B$ et $C$ du triangle : 
\begin{enumerate}
\item Les angles $\widehat{ADB}$ et $\widehat{ADC}$. 
\item Les angles $\widehat{ABE}$ et $\widehat{EBD}$.\\
Application numérique : $\widehat{B}= 68^o$ et $\widehat{C}= 54^o$.
\end{enumerate}
\it Dans un triangle $ABC$, la bissectrice de l'angle $B$ coupe en $I$ la hauteur issue de $A$ et en $D$ la perpendiculaire en $A$ à $(AB)$. 
\begin{enumerate}
\item Évaluer en fonction de $\widehat{B}$ les angles $\widehat{IAB}$, $\widehat{AID}$ et $\widehat{ADI}$.
\item Comparer les segments $[AI]$ et $[AD]$.
\end{enumerate}
\it On considère un cercle $O$ et un diamètre $(xy)$ de ce cercle. D'un point $M$ de ce cercle tel que $\widehat{yOM}<45^o$, on mène $(MH)$ perpendiculaire à $(xy)$ et on construit le point $A$ de $(xy)$ tel que $HA=OH$. La droite $(AM)$ recoupe le cercle en $B$. \begin{enumerate}
\item Montrer que les triangles $MOA$ et $OMB$ sont isocèles. 
\item Calculer les angles $\widehat{OBM}$ et $\widehat{BOx}$ en fonction de l'angle $\widehat{OAB}$.
\end{enumerate}
\it Soit un triangle $ABC$ rectangle en $A$. On prolonge $[CB]$ d'une longueur $BD=BA$ et sur la perpendiculaire en $C$ à $(BC)$, on porte du côté de $A$ une longueur $CE=CA$. 
\begin{enumerate}
\item Calculer en fonction de $\widehat{B}$ les angles $\widehat{BAD}$, $\widehat{ACE}$ et $\widehat{CAE}$. 
\item Évaluer l'angle $\widehat{DAE}$ et montrer que $D$, $A$ et $E$ sont alignés.
\end{enumerate}
\it Soit un triangle rectangle isocèle $ABC$. Par le sommet $A$ de l'angle droit, on mène, extérieurement au triangle une droite $(xy)$, puis les perpendiculaires $(BM)$ et $(CN)$ à $(xy)$ ainsi que la hauteur $(AH)$ du triangle $ABC$. \begin{enumerate}
\item Montrer que $HA=HB=HC$.
\item Comparer les triangles $AMBC$ et $CNA$. Que représente $MN$ pour $BM$ et $CN$ ? 
\item Comparer les triangles $HBM$ et $HAN$ et démontrer que le triangle $HMN$ est rectangle isocèle.
\end{enumerate}
\it On mène la hauteur $(AH)$ issue du sommet de l'angle droit d'un triangle rectangle $ABC$ puis les bissectrices intérieures des angles $\widehat{BAH}$ et $\widehat{CAH}$ qui coupent l'hypoténuse en $D$ et $E$. \begin{enumerate}
\item Évaluer la valeur de l'angle $\widehat{DAE}$.
\item Monter que les triangles $BAE$ et $CAD$ sont isocèles.
\item Comparer $DE$ à la longueur $AB+AC-BC$.
\end{enumerate}
\it Soit un triangle $ABC$. On prolonge $[BC]$ de deux longueurs $BD=BA$ et $CE=CA$. 
\begin{enumerate}
\item Que représente $DE$ pour le triangle $ABC$ ? 
\item Calculer les angles $\widehat{D}$ et $\widehat{E}$ du triangle $ADE$ en fonction des angles $\widehat{B}$ et $\widehat{C}$ de $ABC$.
\item Comparer deux triangles ayant leurs angles respectivement égaux et même périmètre. \end{enumerate}
\it Dans un triangle $ABC$, l'angle aigu $\widehat{B}$ est le double de l'angle $\widehat{C}$. La médiatrice de $[AC]$ coupe $(BC)$ en $D$. 
\begin{enumerate}
\item Montrer que $(AD)$ partage $ABC$ en deux triangles isocèles.
\item Comparer l'angle extérieur au sommet $A$ à l'angle au sommet $C$ dans le triangle $ABC$.
\end{enumerate}
\it Dans un triangle $ABC$, l'angle $\widehat{B}$ est le triple de l'angle $\widehat{C}$. La médiatrice de $[BC]$ coupe $(CA)$ en $D$. 
\begin{enumerate}
\item Montrer que $(BD)$ partage $ABC$ en deux triangles isocèles.
\item Comparer l'angle extérieur au sommet $A$ à l'angle au sommet $C$ dans le triangle $ABC$.
\end{enumerate}
\end{itemize}
